\textbf{BCS teorie:} při nízké teplotě se elektrony s opačnými spiny párují do \textbf{Cooperových párů}. Přitažlivá interakce umožní párování navzdory elektrostatickému odpuzování. Páry vytvářejí společný kvantový koherentní stav, \textbf{BCS kondenzát}, spojený se supravodivostí a nulovou rezistivitou.

Přednáška BCS vysvětluje v souvislosti se \textbf{supravodiči I. třídy}.

Supravodivost zaniká při dosažení \textbf{kritické teploty $T_k$} nebo \textbf{kritické intenzity magnetického pole $H_k$}: Cooperovy páry se rozpadají. Hodnota kritického pole závisí na teplotě, struktuře, čistotě a vnitřním pnutí materiálu.
